\section{隐式 3D 表示}

\subsection{隐式函数的基本思想}

隐式表示通过一个函数 $f(x, y, z) = 0$ 来定义物体表面：满足等式的点在表面上，$f < 0$ 的点在内部，$f > 0$ 的点在外部。

\begin{figure}[H]
\centering
\includegraphics[width=0.95\textwidth]{slides-images/slide-010.jpg}
\caption{隐式表示：单位球面满足 $x^2 + y^2 + z^2 - 1 = 0$。查询点代入函数即可判断内外\protect\footnotemark}
\end{figure}
\footnotetext{Slides 第10页。}

$$\text{单位球：} \quad f(x, y, z) = x^2 + y^2 + z^2 - 1 = 0$$

\begin{importantbox}{隐式表示与显式表示的互补关系}
\begin{itemize}
    \item \textbf{显式表示}：采样容易（直接得到表面点），内外判断困难
    \item \textbf{隐式表示}：内外判断容易（代入函数看符号），采样困难（需要求解方程）
\end{itemize}
这一对偶关系贯穿整个 3D 视觉研究。
\end{importantbox}

\subsection{隐式表示的组合性}

隐式表示的一大优势是\textbf{组合性}（Composability）——可以通过布尔运算组合简单形状来构建复杂形状。

\begin{figure}[H]
\centering
\includegraphics[width=0.95\textwidth]{slides-images/slide-011.jpg}
\caption{隐式函数的布尔运算：并集、交集、差集。通过对函数值做 min/max 运算实现\protect\footnotemark}
\end{figure}
\footnotetext{Slides 第11页。}

给定两个隐式函数 $f_1$ 和 $f_2$：
\begin{itemize}
    \item \textbf{并集}：$\min(f_1, f_2)$
    \item \textbf{交集}：$\max(f_1, f_2)$
    \item \textbf{差集}：$\max(f_1, -f_2)$
\end{itemize}

这种组合能力广泛应用于\textbf{CAD（计算机辅助设计）}和工业制造中。

\begin{knowledgebox}{有符号距离函数（SDF）}
如果隐式函数的值表示到表面的\textbf{有符号距离}（内部为负，外部为正），则称为\textbf{有符号距离函数}（Signed Distance Function, SDF）。SDF 不仅能判断内外，还能通过加法实现形状的\textbf{平滑混合}（Smooth Blending），产生自然的过渡效果。
\end{knowledgebox}

\begin{figure}[H]
\centering
\includegraphics[width=0.95\textwidth]{slides-images/slide-013.jpg}
\caption{SDF 的混合操作：两个距离函数相加即可实现平滑插值。通过组合 SDF 可以构建非常复杂的场景\protect\footnotemark}
\end{figure}
\footnotetext{Slides 第13页。}

\subsection{水平集方法与体素}

对于复杂形状，闭合形式的隐式函数可能不存在。一个实用的方法是：在 3D 网格上\textbf{预计算}隐式函数值，存储为 3D 矩阵——这就是\textbf{水平集方法}（Level Set Method）。

\begin{figure}[H]
\centering
\includegraphics[width=0.95\textwidth]{slides-images/slide-014.jpg}
\caption{水平集方法：在规则网格上预计算距离函数值，相邻值符号变化处即为物体表面\protect\footnotemark}
\end{figure}
\footnotetext{Slides 第14页。}

进一步简化——如果只关心每个网格点是在物体内部还是外部（二值化），就得到了\textbf{体素}（Voxel）表示。

\begin{figure}[H]
\centering
\includegraphics[width=0.95\textwidth]{slides-images/slide-015.jpg}
\caption{体素表示：3D 网格的二值化版本（1=内部, 0=外部），可以看作 3D 版本的像素\protect\footnotemark}
\end{figure}
\footnotetext{Slides 第15页。}

\begin{importantbox}{体素 = 3D 像素}
体素（Voxel）是 2D 像素在 3D 空间的直接推广。像素是 $n \times n$ 矩阵，体素是 $n \times n \times n$ 3D 矩阵。这种类比使得体素成为\textbf{最早被深度学习采用}的 3D 表示——只需将 2D 卷积替换为 3D 卷积即可。
\end{importantbox}

\subsection{本章小结}

\begin{itemize}
    \item 隐式表示通过函数 $f(x,y,z)=0$ 定义物体表面
    \item 优势：内外判断容易，支持布尔组合，SDF 允许平滑混合
    \item 劣势：采样困难，复杂形状难以写出闭合函数
    \item 水平集和体素将隐式函数离散化为 3D 矩阵，便于存储和计算
    \item 体素是最早与深度学习集成的 3D 表示
\end{itemize}

%% ========== 3D 数据集 ==========

